\documentclass[10pt,a4paper]{article}
%\usepackage[utf8x]{inputenc}

%\usepackage[latin1]{inputenc}
%\usepackage[spanish]{babel}

\usepackage{amssymb,latexsym,amsthm,amsmath}
\usepackage{verbatim}
\usepackage{graphicx,color,wrapfig}
\usepackage{psfrag}
\usepackage[all]{xy}
\usepackage{tikz}
%\includepackage{mathabx}
\usepackage{hyperref}
\synctex=1
\usepackage{srcltx}
\hypersetup{colorlinks=true,linkcolor=blue}

\sloppy
\usepackage{multicol}
\setlength{\topmargin}{-0.99in}
\setlength{\textheight}{9.95in}
\setlength{\oddsidemargin}{-0.3in}
\setlength{\evensidemargin}{-0.3in}
\setlength{\textwidth}{6.8in}

\def\R{\mathbb R}
\def\Q{\mathbb Q}
\def\N{\mathbb N}

\DeclareMathOperator{\dom}{dom}
\DeclareMathOperator{\ran}{ran}

\definecolor{light-gray}{gray}{0.85}


\usepackage{type1cm}
\usepackage{eso-pic}
\usepackage{calc}
\usepackage{forloop}

\newcounter{ber}
\newcounter{num}

\newcommand{\lef}{\langle}
\newcommand{\rig}{\rangle}

\makeatletter

\makeatother



\begin{document}

\begin{center}
\begin{figure}[ht!]
\begin{minipage}[c]{0.2\linewidth}
\includegraphics[scale=0.22]{EscudoUNI.eps}
\end{minipage}
\hfill{}
\begin{minipage}{0,8\columnwidth}
\centering\flushleft {\large\bf Introducci\'on a la Programaci\'on\\
Escuela Profesional de Ciencia de la Computaci\'on\\
Facultad de Ciencias\\
Universidad Nacional de Ingenier\'i�a\\}
\end{minipage}
\end{figure}
\end{center}
\vspace*{-1.2cm}
\noindent

\begin{center}
\textbf{\Large Laboratorio dirigido 4}
\end{center}

\textbf{Tema: Funciones en C \hfill Ciclo: 2015-I}
\noindent
\rule{\textwidth}{1pt}
\newline
\vspace*{-0.5cm}



\begin{enumerate}

\item Elaborar un programa que imprima el men\'u mostrado abajo, pida la cantidad a convertir seg\'un la opci\'on ingresada y muestre la conversi\'on. Para ello, se deber\'a crear funciones para realizar cada conversi\'on.
 \begin{itemize}
 \item[1.] Cent\'igrados a Fahrenheit
 \item[2.] Fahrenheit a cent\'igrados
 \item[3.] Kilogramos a Libras
 \item[4.] Libras a Kilogramos
 \item[5.] Salir 
 \end{itemize}
\item Cree un programa que pida al usuario ingresar un entero positivo $n$ y muestre, llamando a una funci\'on primo({\bf int} $n$) que retorna 1 si se ingresa un primo y 0 si se ingresa un entero que no es primo, todos los primos menores o iguales a $n$.
\item Haga un programa que nos permita introducir un entero positivo $n$ y sobre \'el se realicen las siguientes operaciones: comprobar si es primo, hallar su factorial y hallar la suma desde 1 hasta $n$. Cada operaci\'on debe realizarse en una funci\'on. El men\'u debe aparecer hasta que el usuario indique que ha finalizado.
\item Implemente un program que pida ingresar dos enteros positivos y muestre su MCD y MCM evocando dos funciones mcd({\bf int} a, {\bf int} b) y mcm({\bf int} a, {\bf int} b) que retornan el MCD y MCM, respectivamente, si se ingresan dos enteros positivos.
%\item Elabore un programa que pida ingresar tres enteros y muestre el m\'inimo, m\'aximo y la media aritm\'etica de ellos mediante tres funciones: min({\bf int} a, {\bf int} b, {\bf int} c), max({\bf int} a, {\bf int} b, {\bf int} c) y med({\bf int} a, {\bf int} b, {\bf int} c).
\item Cree un programa que lea dos vectores de 3 componentes. Luego, los sume, reste y multiplique, tanto escalar como vectorialmente, imprimiendo después todos estos resultados. Para ello se deber\'a implementar las funciones suma(), resta(), productoEscalar() y productoVectorial().
\item Implemente un programa que pida el valor de los coeficientes $a$, $b$ y $c$ de la ecuaci\'on cuadr\'atica $$ ax^2 + bx + c = 0,\qquad\mbox{donde}\quad a,b,c\in\R, $$ y muestre las ra\'ices de dicha ecuaci\'on. Para ello se deber\'a implementar tres funciones: realesDiferentes({\bf float} a, {\bf float} b, {\bf float} Delta), realesIguales({\bf float} a, {\bf float} b) y complejas({\bf float} a, {\bf float} b, {\bf float} Delta) que mostrar\'an las ra\'ices en caso el discriminante sea positivo, cero o negativo, respectivamente. 
\item Implemente un programa tal que al ingresar un entero $n>1$ dibuje:
 \begin{itemize}
 \item un rombo cuyas diagonales miden $2n+1$.
 \item un tri\'angulo de altura $2n-1$ y ancho $n$.
 \item un tri\'angulo de altura $n$ y ancho $2n-1$.
 \end{itemize}
 Para ello se debe implementar 3 funciones: rombo({\bf int} n), tri1({\bf int} n) y tri2({\bf int} n). Por ejemplo, para $n=2$ se tendr\'a:\\
 \begin{center}
 \begin{tabular}{ccccc}
 &&*&&\\
 &*&*&*&\\
 *&*&*&*&*\\
 &*&*&*&\\
 &&*&&\\ 
 \\
 *&&&&\\
 *&*&&&\\
 *&&&&\\
 \\
  &*&&&\\
 *&*&*&&\\
 \end{tabular}
 \end{center}
 
\end{enumerate}

\end{document}
